\documentclass[12pt,a4paper]{scrartcl}
\usepackage[utf8]{inputenc}
%\usepackage[latin1]{inputenc} %  Alternativ unter Windows
\usepackage[T1]{fontenc}
\usepackage[paper=a4paper,
	left=25mm,
	right=25mm,
	top=25mm,
	bottom=20mm]{geometry}
\usepackage[nottoc]{tocbibind}
\usepackage[ngerman]{babel}
\usepackage[pdftex]{graphicx}
\usepackage{latexsym}
\usepackage{amsmath,amssymb,amsthm}
\usepackage{mathtools}

\newtheorem{Satz}{Satz}[section]
\newtheorem{Definition}[Satz]{Definition} 
\newtheorem{Lemma}[Satz]{Lemma}		                 
\newtheorem{Korollar}[Satz]{Korollar}
\newtheorem{Proposition}[Satz]{Proposition}                  
\newtheorem{examples}[Satz]{Beispiel}
\let\oldexamples\examples
\renewcommand{\examples}{\oldexamples\normalfont}
\newtheorem*{notation}{Notation}
\let\oldnotation\notation
\renewcommand{\notation}{\oldnotation\normalfont}
\newtheorem*{remark}{Bemerkung}
\let\oldremark\remark
\renewcommand{\remark}{\oldremark\normalfont}                  
                              
\numberwithin{equation}{section} 

% Praktikumsbericht nummer setzen
\newcommand{\BerichtNR}{1}

\begin{document}
\begin{titlepage}
	\includegraphics[scale=0.5]{kit-logo.jpg} 
	\begin{center} 
		\LARGE 
		\vspace*{2cm}
		\LARGE Praktikumsbericht \BerichtNR
		\vspace*{1.0cm}
		\hrule
		\vspace*{0.2cm}
		{\vspace{0.2cm} \huge Einführung in das Wissenschaftliche Rechnen}\vspace{0.5cm}
		\hrule
		\vspace*{2.5cm}
		\Large Stefan Karch \\
		Florian Döttling \\
		Tim Buchholz  \\
		\vspace*{1cm}
		20.05.2019 \\
		\vspace*{1.5cm}
		\vspace*{4.0cm}
		\Large Betreuung: Prof. Dr. Christian Wieners, Niklas Baumgarten \\[0.5cm]
		\Large Fakultät für Mathematik \\
		\Large Karlsruher Institut für Technologie
	\end{center}
\end{titlepage}

%\tableofcontents
%\newpage 
%\pagestyle{headings}

\section{Problemstellung}

TODO: Beschreibung der Problemstellung, Beschreibung der Aufgabenstellung und Beschreibung der numerischen Verfahren


\section{Ergebnisse}

TODO: Dokumentation der im Praktikum erzielten Ergebnisse 


\end{document}

